% GENERATED FILE. Edit canonical lessons, then run npm run generate:books.
% canonical-source-hash: fnv1a64:4e81a9b320d9b470
% canonical-lessons: JA-C360-shuunyuu, JA-C360-shishutsu, JA-C360-seikatsuhi, JA-C360-tamaru, JA-C360-osaeru

\chapter{Income and Spending}
\label{ch:ja-income-and-spending}
\hlchaptermodality{360}

\begin{chapteropening}
\textbf{By the end of this chapter:} \emph{I can say what I would do if my income or spending changed.}

Complete the last lesson of chapter 360: I can say what I would do if my income or spending changed.
\end{chapteropening}

\section[\texorpdfstring{\ja{しゅうにゅう} (shūnyū)}{しゅうにゅう (shūnyū)}]{\ja{しゅうにゅう} (shūnyū) — income}

\label{lesson:JA-C360-shuunyuu}



\begin{quote}
Before the new one: say the Japanese for a refund, then the Japanese for a sum of money.
\end{quote}

\subsection*{You'll want to know: \ja{しゅうにゅう}}
\textbf{\ja{しゅうにゅう}} — \emph{shūnyū} — \textquotedblleft{}income\textquotedblright{}.

\begin{quote}
\textbf{\ja{しゅうにゅうが} \ja{ふえたら、ひっこしたいです。}} — \emph{shūnyū ga fuetara, hikkoshitai desu} — \textquotedblleft{}If my income goes up, I want to move house.\textquotedblright{}
\end{quote}

\begin{cousinweb}
\textbf{\ja{しゅうにゅう}} (\emph{shūnyū}) is Sino-Japanese (kango): two Chinese characters read with their on'yomi, readings borrowed from Chinese. They are \emph{shū}, \textquotedblleft{}income\textquotedblright{} (Middle Chinese \emph{syuw}) and \emph{nyū}, \textquotedblleft{}enter\textquotedblright{} (\emph{nyip}) (KANJIDIC2; Middle Chinese in Baxter's transcription).
\end{cousinweb}

\begin{grammarlens}[title={What you've built}]
One more word for saying what happens if.
\end{grammarlens}

\subsection*{Your turn}
\begin{itemize}
\raggedright
  \item \emph{Say it:} \emph{shūnyū}
  \item \emph{Say it:} \emph{shūnyū}, once more
  \item \emph{Say it:} \emph{shūnyū ga fuetara, hikkoshitai desu} — If my income goes up, I want to move house.
  \item \emph{Recall:} say the Japanese for a refund, then the Japanese for a sum of money, then say \emph{shūnyū} again
\end{itemize}

\begin{tcolorbox}[breakable,colback=teal!4,colframe=teal!35!black,title={Before you move on}]
What does \ja{しゅうにゅう} mean? (\textquotedblleft{}Income\textquotedblright{}.) Say it once more.
\end{tcolorbox}
\section[\texorpdfstring{\ja{ししゅつ} (shishutsu)}{ししゅつ (shishutsu)}]{\ja{ししゅつ} (shishutsu) — spending; expenditure}

\label{lesson:JA-C360-shishutsu}



\begin{quote}
Before the new one: say the Japanese for a sum of money, then the Japanese for income.
\end{quote}

\subsection*{You'll want to know: \ja{ししゅつ}}
\textbf{\ja{ししゅつ}} — \emph{shishutsu} — \textquotedblleft{}spending; expenditure\textquotedblright{}.

\begin{quote}
\textbf{\ja{ししゅつが} \ja{ふえたら、せつやくした} \ja{ほうが} \ja{いいです。}} — \emph{shishutsu ga fuetara, setsuyaku shita hō ga ii desu} — \textquotedblleft{}If spending goes up, we had better economize.\textquotedblright{}
\end{quote}

\begin{cousinweb}
\textbf{\ja{ししゅつ}} (\emph{shishutsu}) is Sino-Japanese (kango): two Chinese characters read with their on'yomi, readings borrowed from Chinese. They are \emph{shi}, \textquotedblleft{}support\textquotedblright{} (Middle Chinese \emph{tsye}) and \emph{shutsu}, \textquotedblleft{}exit\textquotedblright{} (\emph{tsyhwit}) (KANJIDIC2; Middle Chinese in Baxter's transcription).
\end{cousinweb}

\begin{grammarlens}[title={What you've built}]
One more word for saying what happens if.
\end{grammarlens}

\subsection*{Your turn}
\begin{itemize}
\raggedright
  \item \emph{Say it:} \emph{shishutsu}
  \item \emph{Say it:} \emph{shishutsu}, once more
  \item \emph{Say it:} \emph{shishutsu ga fuetara, setsuyaku shita hō ga ii desu} — If spending goes up, we had better economize.
  \item \emph{Recall:} say the Japanese for a sum of money, then the Japanese for income, then say \emph{shishutsu} again
\end{itemize}

\begin{tcolorbox}[breakable,colback=teal!4,colframe=teal!35!black,title={Before you move on}]
What does \ja{ししゅつ} mean? (\textquotedblleft{}Spending; expenditure\textquotedblright{}.) Say it once more.
\end{tcolorbox}
\section[\texorpdfstring{\ja{せいかつひ} (seikatsuhi)}{せいかつひ (seikatsuhi)}]{\ja{せいかつひ} (seikatsuhi) — living costs}

\label{lesson:JA-C360-seikatsuhi}



\begin{quote}
Before the new one: say the Japanese for income, then the Japanese for expenditure.
\end{quote}

\subsection*{You'll want to know: \ja{せいかつひ}}
\textbf{\ja{せいかつひ}} — \emph{seikatsuhi} — \textquotedblleft{}living costs\textquotedblright{}.

\begin{quote}
\textbf{\ja{とかいに} \ja{すんだら、せいかつひが} \ja{たかいでしょう。}} — \emph{tokai ni sundara, seikatsuhi ga takai deshō} — \textquotedblleft{}If you live in the city, your living costs will probably be high.\textquotedblright{}
\end{quote}

\begin{cousinweb}
\textbf{\ja{せいかつひ}} (\emph{seikatsuhi}) is Sino-Japanese (kango): three Chinese characters read with their on'yomi, readings borrowed from Chinese. They are \emph{sei}, \textquotedblleft{}life\textquotedblright{} (Middle Chinese \emph{srjaeng}), \emph{katsu}, \textquotedblleft{}lively\textquotedblright{} (\emph{hwat}) and \emph{hi}, \textquotedblleft{}expense\textquotedblright{} (\emph{phj+jH}) (KANJIDIC2; Middle Chinese in Baxter's transcription).
\end{cousinweb}

\begin{grammarlens}[title={What you've built}]
One more word for saying what happens if.
\end{grammarlens}

\subsection*{Your turn}
\begin{itemize}
\raggedright
  \item \emph{Say it:} \emph{seikatsuhi}
  \item \emph{Say it:} \emph{seikatsuhi}, once more
  \item \emph{Say it:} \emph{tokai ni sundara, seikatsuhi ga takai deshō} — If you live in the city, your living costs will probably be high.
  \item \emph{Recall:} say the Japanese for income, then the Japanese for expenditure, then say \emph{seikatsuhi} again
\end{itemize}

\begin{tcolorbox}[breakable,colback=teal!4,colframe=teal!35!black,title={Before you move on}]
What does \ja{せいかつひ} mean? (\textquotedblleft{}Living costs\textquotedblright{}.) Say it once more.
\end{tcolorbox}
\section[\texorpdfstring{\ja{たまる} (tamaru)}{たまる (tamaru)}]{\ja{たまる} (tamaru) — to be saved up; to build up}

\label{lesson:JA-C360-tamaru}



\begin{quote}
Before the new one: say the Japanese for expenditure, then the Japanese for living costs.
\end{quote}

\subsection*{You'll want to know: \ja{たまる}}
\textbf{\ja{たまる}} — \emph{tamaru} — \textquotedblleft{}to be saved up; to build up\textquotedblright{}.

\begin{quote}
\textbf{\ja{おかねが} \ja{たまったら、りょこうに} \ja{いきたいです。}} — \emph{okane ga tamattara, ryokō ni ikitai desu} — \textquotedblleft{}When I have saved enough money, I want to go travelling.\textquotedblright{}
\end{quote}

Polite \textbf{\ja{たまります}}, \emph{tamarimasu}; a godan verb. Past (polite): \textbf{\ja{ごみが} \ja{たまりました。}}, \emph{gomi ga tamarimashita}.

\begin{cousinweb}
\textbf{\ja{たまる}} (\emph{tamaru}) is a native Japanese word (wago), inherited rather than borrowed from Chinese, and its older origin is not settled. JMdict notes that it is usually written in kana (JMdict).
\end{cousinweb}

\begin{grammarlens}[title={What you've built}]
One more word for saying what happens if.
\end{grammarlens}

\subsection*{Your turn}
\begin{itemize}
\raggedright
  \item \emph{Say it:} \emph{tamaru}
  \item \emph{Say it:} \emph{tamaru}, once more
  \item \emph{Say it:} \emph{okane ga tamattara, ryokō ni ikitai desu} — When I have saved enough money, I want to go travelling.
  \item \emph{Recall:} say the Japanese for expenditure, then the Japanese for living costs, then say \emph{tamaru} again
\end{itemize}

\begin{tcolorbox}[breakable,colback=teal!4,colframe=teal!35!black,title={Before you move on}]
What does \ja{たまる} mean? (\textquotedblleft{}To be saved up; to build up\textquotedblright{}.) Say it once more.
\end{tcolorbox}
\section[\texorpdfstring{\ja{おさえる} (osaeru)}{おさえる (osaeru)}]{\ja{おさえる} (osaeru) — to keep down (costs); to hold down}

\label{lesson:JA-C360-osaeru}



\begin{quote}
Before the new one: say the Japanese for living costs, then the Japanese for to be saved up.
\end{quote}

\subsection*{You'll want to know: \ja{おさえる}}
\textbf{\ja{おさえる}} — \emph{osaeru} — \textquotedblleft{}to keep down (costs); to hold down\textquotedblright{}.

\begin{quote}
\textbf{\ja{ねだんが} \ja{あがったら、ししゅつを} \ja{おさえた} \ja{ほうが} \ja{いいです。}} — \emph{nedan ga agattara, shishutsu o osaeta hō ga ii desu} — \textquotedblleft{}If prices go up, we had better keep our spending down.\textquotedblright{}
\end{quote}

Polite \textbf{\ja{おさえます}}, \emph{osaemasu}; an ichidan verb. Past (polite): \textbf{\ja{ははが} \ja{ししゅつを} \ja{おさえました。}}, \emph{haha ga shishutsu o osaemashita}.

\begin{cousinweb}
\textbf{\ja{おさえる}} (\emph{osaeru}) is a native Japanese word (wago), inherited rather than borrowed from Chinese, and its older origin is not settled. Its kanji was attached to it later as a kun'yomi, a native reading given to a Chinese character (JMdict, KANJIDIC2).
\end{cousinweb}

\begin{grammarlens}[title={What you've built}]
One more word for saying what happens if.
\end{grammarlens}

\subsection*{Your turn}
\begin{itemize}
\raggedright
  \item \emph{Say it:} \emph{osaeru}
  \item \emph{Say it:} \emph{osaeru}, once more
  \item \emph{Say it:} \emph{nedan ga agattara, shishutsu o osaeta hō ga ii desu} — If prices go up, we had better keep our spending down.
  \item \emph{Recall:} say the Japanese for living costs, then the Japanese for to be saved up, then say \emph{osaeru} again
\end{itemize}

\begin{tcolorbox}[breakable,colback=teal!4,colframe=teal!35!black,title={Before you move on}]
What does \ja{おさえる} mean? (\textquotedblleft{}To keep down (costs); to hold down\textquotedblright{}.) Say it once more.
\end{tcolorbox}
